\chapter{Lab Solutions and Solution Sketches}
\index{lab solutions}\index{hands-on projects!solutions}%
Solution sketches for every chapter's Thursday hands-on project (Chapters 1--24) and selected
Chapter 0 exercises. Each entry gives the intended approach, the key commands or code shape, and
the failure modes that cost students the most time. These are sketches, not the only correct
answers: if your solution meets the chapter's acceptance criteria with different code, it counts.

\section{Chapter 0: Setup Verification Exercises}
\textbf{Approach.} Complete both paths---Atlas M0 and local Community Server---and capture the
verification transcript for each: ping, insert, find, and (locally)
\texttt{db.serverStatus().storageEngine} showing \texttt{wiredTiger}.

\textbf{Key commands.}
\begin{lstlisting}[language=bash, caption={Chapter 0 verification, both environments.}]
mongosh "mongodb+srv://bookUser:<password>@cluster0.example.mongodb.net/booklabs"
mongosh "mongodb://localhost:27017/booklabs"
# inside mongosh, both targets:
db.runCommand({ ping: 1 })
db.chapter0.insertOne({ environment: "atlas" })   # or "local"
db.chapter0.findOne({ environment: "atlas" })
db.serverStatus().storageEngine                    # local only
\end{lstlisting}

\textbf{Common failure modes.} Authentication failure = wrong database user or an unencoded
reserved character in the password; timeout = the IP access list no longer matches your current
public IP; \texttt{mongodb+srv} DNS failure = blocked DNS or VPN intercept; local service down =
read \texttt{mongod.cfg}'s log path before touching anything else.

\section{Chapter 1: Configure and Observe WiredTiger}
\textbf{Approach.} Run a local \texttt{mongod} with a development-only, explicitly small
WiredTiger cache, generate a write workload larger than the cache, and watch
\texttt{serverStatus().wiredTiger.cache} metrics move: bytes in cache, dirty bytes, eviction, and
checkpoint counters.

\textbf{Key commands.} Set \texttt{storage.wiredTiger.engineConfig.cacheSizeGB} in a lab
\texttt{mongod.conf}; drive writes with a \texttt{mongosh} loop of
\texttt{insertMany} batches; poll \texttt{db.serverStatus().wiredTiger.cache} and
\texttt{.concurrentTransactions} between batches.

\textbf{Common failure modes.} Editing the production service's \texttt{mongod.cfg} instead of a
lab config; cache too small to start (keep $\geq$ 0.25 GB); reading averages instead of p99 and
concluding ``no checkpoint effect''; forgetting that the OS page cache masks small datasets.

\section{Chapter 2: Refactor a Product Catalog}
\textbf{Approach.} Take the monolithic product document (sparse dynamic fields, unbounded
reviews array) and split it: Attribute Pattern key-value array for specifications, Subset Pattern
for the ten hottest reviews embedded, full reviews in a child collection. Migrate with a PyMongo
script that writes the new shape and validates counts.

\textbf{Key code shape.} \texttt{createIndex(\{"attrs.k": 1, "attrs.v": 1\})} for dynamic filters;
migration loops \texttt{find()} batches, builds \texttt{bulkWrite} operations, and asserts
\texttt{countDocuments} parity per source.

\textbf{Common failure modes.} Migrating without a schema-version field so readers cannot tell
old from new; embedding \emph{all} reviews (the unbounded-array bug returns); indexing each
sparse spec field separately instead of the Attribute Pattern pair; no rollback script.

\section{Chapter 3: Enterprise HR Validator}
\textbf{Approach.} Create an \texttt{employees} collection with a \texttt{\$jsonSchema}
validator covering required fields, types, enums, nested shapes, and array bounds; then insert
valid and invalid documents from PyMongo and assert the rejection paths.

\textbf{Key code shape.}
\begin{lstlisting}[language=JavaScript, caption={Validator skeleton with strict rejection.}]
db.createCollection("employees", { validator: { $jsonSchema: {
  bsonType: "object", required: ["employeeId", "email", "roles"],
  properties: {
    email: { bsonType: "string", pattern: "^[^@]+@[^@]+$" },
    roles: { bsonType: "array", maxItems: 5,
             items: { enum: ["admin","editor","viewer"] } } } } },
  validationLevel: "strict", validationAction: "error" })
\end{lstlisting}

\textbf{Common failure modes.} \texttt{validationAction: "warn"} left on, so bad documents
silently land; pattern regexes that reject valid emails; testing only the happy path; forgetting
that validators do not apply to pre-existing documents unless writes touch them.

\section{Chapter 4: E-Commerce Relational Migrator Lab}
\textbf{Approach.} Map twelve PostgreSQL tables into three MongoDB collections---orders with
embedded line items, customers with embedded addresses, products---using Relational Migrator (or
the provided mapping YAML); run initial sync, then CDC, then verify counts, document shape, and
change propagation with a PyMongo script.

\textbf{Key verification.} Per-table count reconciliation; spot-check embedded array cardinality
against SQL \texttt{COUNT} grouped by parent; update a source row and time its appearance in Atlas.

\textbf{Common failure modes.} One-table-to-one-collection mapping (the anti-pattern the chapter
warns against); verifying counts while CDC is behind and declaring false divergence; cutover
planned without a rollback authority (which system owns writes during reversal).

\section{Chapter 5: Break a Replica Set on Purpose}
\textbf{Approach.} Deploy three local \texttt{mongod} instances (ports 28017--28019) as
\texttt{labrs}, then attack it: \texttt{rs.stepDown()}, kill a secondary, and partition the
network with Windows firewall rules, writing continuously at \texttt{w:majority} throughout and
recording availability windows.

\textbf{Key commands.} \texttt{rs.initiate} with three members; a writer loop catching errors;
\texttt{rs.status()} to watch term changes; \texttt{rs.printReplicationInfo()} for the oplog
window; \texttt{New-NetFirewallRule -Block} to simulate the partition.

\textbf{Common failure modes.} Forgetting \texttt{--oplogSize} so the default dwarfs the lab
disk; reading \texttt{rs.status()} too late and missing the election; firewall rules left behind
after the lab; concluding \texttt{w:1} is ``fine'' because rollback did not occur in a two-minute
test.

\section{Chapter 6: A Two-Shard Cluster on One Machine}
\textbf{Approach.} Build the full topology locally: config server replica set, two shard
\texttt{mongod} processes, and one \texttt{mongos}; enable sharding on a database, shard a
collection with a hashed key, insert bulk data, and watch chunks form and distribute via
\texttt{sh.status()} and the \texttt{config.chunks} collection.

\textbf{Key commands.} \texttt{mongod --configsvr --replSet cfgrs}; \texttt{mongod --shardsvr};
\texttt{mongos --configdb cfgrs/...}; then \texttt{sh.addShard(...)},
\texttt{sh.enableSharding("booklabs")}, \texttt{sh.shardCollection("booklabs.events", \{userId: "hashed"\})}.

\textbf{Common failure modes.} Starting \texttt{mongos} before the config server is initialized;
querying a shard directly and concluding data ``disappeared'' (always query through
\texttt{mongos}); undersized disks when four \texttt{mongod} processes share one drive.

\section{Chapter 7: Break a Hot Shard, Then Refine It Live}
\textbf{Approach.} Shard a collection on a low-cardinality region key, drive writes at one
region until one shard is measurably hot and a jumbo chunk forms, then run
\texttt{refineCollectionShardKey} to append a bucket suffix \emph{without downtime} and watch the
balancer restore divisibility.

\textbf{Key commands.} Generate documents with a \texttt{bucket} field cycling 0--63;
\texttt{db.adminCommand(\{refineCollectionShardKey: "db.coll", key: \{region:1, \_id:1, bucket:1\}\})};
observe with \texttt{sh.status()} and per-shard \texttt{serverStatus().opcounters}.

\textbf{Common failure modes.} Attempting to change the key \emph{prefix} (requires
\texttt{reshardCollection} or migration, not refinement); not backfilling the suffix field on
existing documents; measuring ``hot'' by latency alone instead of shards-hit and per-shard ops.

\section{Chapter 8: Zone-Pinned User Residency Lab}
\textbf{Approach.} Tag shards as \texttt{EU} and \texttt{US}, map shard-key ranges of the
country-code prefix to each zone, insert users with EU and US codes, and verify residency by
reading the chunk map---not by trusting the insert path.

\textbf{Key commands.} \texttt{sh.addShardToZone("shard0","EU")};
\texttt{sh.updateZoneKeyRange("shop.users", \{cc:"A"\}, \{cc:"M"\}, "EU")}; verification loops
\texttt{config.chunks} and asserts every EU-intersecting chunk lives on an EU-zoned shard.

\textbf{Common failure modes.} Country codes not covered by any zone range floating anywhere
(the silent violation); residency field not leading the shard key, so chunks straddle
jurisdictions; proving residency once by hand instead of scheduling the assertion.

\section{Chapter 9: Prove the ESR Rule on Five Million Documents}
\textbf{Approach.} Generate five million orders, then benchmark the same query under three
plans: no index, an ESR-violating compound index (range or sort first), and the correct
equality--sort--range index, comparing \texttt{totalDocsExamined}, \texttt{nReturned}, and
execution time. Finish with a rolling index build rehearsal.

\textbf{Key commands.} \texttt{db.orders.explain("executionStats").find(\{status:"A"\}).sort(\{orderDate:-1\}).limit(50)};
\texttt{createIndex(\{status:1, orderDate:-1\})} as the corrected version; rolling build by
stopping one secondary as a standalone, building, rejoining, repeating, and stepping down last.

\textbf{Common failure modes.} Datasets small enough to fit in cache so every plan looks fast;
comparing cold-cache to warm-cache runs; \texttt{SORT\_KEY\_GENERATOR} missed in the plan tree;
rolling build attempted on Atlas M0 (it hides the nodes---use the local set).

\section{Chapter 10: Organizational Hierarchy Analytics}
\textbf{Approach.} Generate a synthetic org chart, compute recursive headcount rollups with
\texttt{\$graphLookup} (bounded by \texttt{maxDepth}, pruned by
\texttt{restrictSearchWithMatch}), then join to budgets with \texttt{\$lookup} and reshape with
\texttt{\$group}/\texttt{\$facet}.

\textbf{Key code shape.} \texttt{\$graphLookup: \{from:"employees", startWith:"\$managerId",
connectFromField:"managerId", connectToField:"\_id", as:"reports", maxDepth:10\}}; foreign side
of every \texttt{\$lookup} indexed.

\textbf{Common failure modes.} Unbounded \texttt{\$graphLookup} on cyclic data; \texttt{\$facet}
output exceeding 16 MB; \texttt{\$match} placed after a reshaping stage, defeating pushdown.

\section{Chapter 11: Financial Moving-Average Pipeline}
\textbf{Approach.} Generate a ledger, then compute moving averages, ranks, and day-over-day
derivatives with \texttt{\$setWindowFields} (\texttt{sortBy} on the date), and persist results
with \texttt{\$merge} using a deterministic key so re-runs are idempotent. Add a lookback window
to the incremental refresh so late events re-compute the frames they touch.

\textbf{Key code shape.} \texttt{\$setWindowFields: \{partitionBy:"\$account", sortBy:\{date:1\},
output:\{ma30:\{\$avg:"\$amount", window:\{range:[-30,"current"], unit:"day"\}\}\}\}}; then
\texttt{\$merge: \{into:"daily\_stats", on:["account","date"], whenMatched:"replace"\}}.

\textbf{Common failure modes.} Missing \texttt{sortBy} (nondeterministic frames);
\texttt{\$rank} where the business wanted \texttt{\$denseRank}; refresh without lookback, so
late-arriving ledger entries silently corrupt history; non-idempotent merge keys duplicating rows.

\section{Chapter 12: The Un-Optimized Database Audit}
\textbf{Approach.} Given a deliberately neglected database (redundant indexes, a spilling
report, fat documents), run the audit sequence: profile to find slow operations, rank by total
cost, \texttt{explain} the top offenders, fix ESR order, hide---then drop---unused indexes found
via \texttt{\$indexStats}, and re-measure.

\textbf{Key commands.} \texttt{db.setProfilingLevel(1, \{slowms:50\})};
\texttt{db.system.profile.aggregate} grouped by query shape with \texttt{\$sum:"\$millis"};
\texttt{db.coll.aggregate([\{\$indexStats:\{\}\}])}; \texttt{db.coll.hideIndex(...)} before any
\texttt{dropIndex}.

\textbf{Common failure modes.} Optimizing the worst-latency query instead of the
highest-total-cost one; dropping an index inside its business cycle (month-end report breaks in
three weeks); auditing only the primary while the analytics load runs on secondaries.

\section{Chapter 13: Autocomplete Product Search}
\textbf{Approach.} Define an Atlas Search index with an \texttt{autocomplete} (edge-ngram)
mapping on the product title plus a token field for category, load the sample catalog, poll index
status until caught up, and serve type-ahead queries through \texttt{\$search} as the first
pipeline stage with facets for filters.

\textbf{Key code shape.} Index JSON per the quick-reference appendix; query with
\texttt{compound: \{must: [autocomplete], filter: [equals on category]\}}; \texttt{\$searchMeta}
for facet counts.

\textbf{Common failure modes.} Serving traffic before the index syncs (incomplete results, no
error); leading-wildcard regex instead of edge-ngram tokenization; \texttt{\$search} not in first
position; ngram fields ballooning index size on long text.

\section{Chapter 14: The RAG Backend}
\textbf{Approach.} Chunk the source documents, embed each chunk with the chosen provider, store
chunks + vectors + tenant metadata in Atlas, define a vector index declaring every filter field,
and answer questions by retrieving with \texttt{\$vectorSearch} (pre-filtered by tenant) and
grounding the LLM prompt in the retrieved chunks.

\textbf{Key code shape.} \texttt{\$vectorSearch: \{index:"vec", path:"embedding",
queryVector:qvec, numCandidates:200, limit:10, filter:\{tenantId:\{\$eq:t\}\}\}}; score via
\texttt{\$meta:"vectorSearchScore"}; hybrid variant fuses BM25 \texttt{\$search} and vector
results with reciprocal rank fusion.

\textbf{Common failure modes.} Filter fields undeclared in the index (over-fetch + cross-tenant
risk); chunk sizes that break sentence semantics; Euclidean distance on text embeddings trained
for cosine; no eval set, so prompt/index changes regress silently.

\section{Chapter 15: Exactly-Once Kafka Windowed Pipeline}
\textbf{Approach.} Wire a Kafka source into Atlas Stream Processing, compute tumbling-window
aggregates with \texttt{\$group} over event-time windows, and sink with \texttt{\$merge} keyed
deterministically on (window start, group key) so checkpoint replays re-emit harmlessly---
at-least-once delivery absorbed by an idempotent sink.

\textbf{Common failure modes.} Non-deterministic merge keys (duplicates on every restart);
unbounded state from a high-cardinality group key; watermark too tight so late events are
dropped silently; processor downtime longer than the oplog window invalidating resume state.

\section{Chapter 16: Search-Index Sync Microservice}
\textbf{Approach.} Consume a change stream, transform events, and upsert into a target
collection or external index. Persist the resume token \emph{after} each successful sink write;
on restart, resume with \texttt{resumeAfter}. Wrap the business write + event in an outbox where
the consumer must not miss non-streamed writes.

\textbf{Key code shape.} \texttt{coll.watch(pipeline, \{resumeAfter: savedToken\})}; token saved
in a \texttt{consumer\_state} collection post-write; idempotent sink keyed on document
\texttt{\_id} + operation time.

\textbf{Common failure modes.} Saving the token before the sink write (event loss on crash);
serverless consumers cold-starting unbounded connection pools; consumer paused past the oplog
window (re-bootstrap from snapshot required); ``write then publish'' without an outbox.

\section{Chapter 17: x.509 Cluster with Least-Privilege Roles}
\textbf{Approach.} Stand up a local TLS-enabled cluster (or Atlas dedicated tier), issue client
certificates from a lab CA with OpenSSL, authenticate services by certificate, and replace broad
roles with custom roles granting per-collection, per-action privileges matching each service's
documented access patterns.

\textbf{Key commands.} \texttt{db.createRole(\{role:"ordersReader", privileges:[\{resource:
\{db:"shop",collection:"orders"\}, actions:["find"]\}], roles:[]\})}; certificate login via
\texttt{--tls --tlsCertificateKeyFile} and \texttt{\$external} users.

\textbf{Common failure modes.} Certificate subject DN mismatch between user record and client
cert; role edited in the shell and never captured in code; \texttt{AnyDatabase} roles creeping in
``temporarily''; no alert on authorization-failure spikes after rollout.

\section{Chapter 18: Queryable Encryption with AWS KMS}
\textbf{Approach.} Configure a KMS customer master key, create data encryption keys per field,
connect PyMongo with the QE shared library and \texttt{AutoEncryptionOpts}, insert a document,
and query an encrypted equality field---verifying in the shell (without keys) that only
ciphertext is stored.

\textbf{Key code shape.} \texttt{ClientEncryption(kms\_providers=\{"aws":\{...\}\},
key\_vault\_namespace="encryption.\_\_keyVault", ...)}; collection created with
\texttt{encryptedFields} mapping field paths to DEK IDs; query with the QE-enabled client returns
plaintext, a plain client returns ciphertext.

\textbf{Common failure modes.} Querying with a client that lacks encryption config and seeing
garbage; deterministic encryption on a low-cardinality field; underestimating the 2--4$\times$
storage cost; no tested procedure for KMS key rotation and (for erasure) destruction.

\section{Chapter 19: Auditing a Sensitive Collection}
\textbf{Approach.} Enable Atlas auditing with a filter scoped to identity events plus DML on one
sensitive collection, generate representative traffic (allowed and denied), ship logs off-host,
and write the detection query that finds anomalous access---then verify the audit pipeline's own
health is monitored.

\textbf{Common failure modes.} Full-fidelity auditing drowning the signal and the disk; audit
logs readable only on the host they describe; filter drift as new sensitive collections appear;
alert rules tested against synthetic events but never against the real pipeline.

\section{Chapter 20: Online Archive with a Federated Query}
\textbf{Approach.} Create an Online Archive rule on a large time-series collection partitioned
by the query date key, let aging data tier to object storage, and run an Atlas Data Federation
query that unions live and archived data---measuring bytes scanned to confirm partition pruning
works.

\textbf{Common failure modes.} Partition key not matching the query predicate (unbounded scans,
unbounded invoice); treating archive lag as replication lag; no documented path for
\textbf{un}-archiving a mistaken rule; TTL and Online Archive rules fighting over the same data.

\section{Chapter 21: The Three-Environment GitOps Pipeline}
\textbf{Approach.} Structure a repository as \texttt{envs/dev}, \texttt{envs/stage},
\texttt{envs/prod} Terraform roots sharing modules; CI runs \texttt{terraform plan} on pull
requests and \texttt{apply} on merge, with remote state locked and Atlas API keys injected from
the CI secret store. Bundle the schema-contract test (validators, indexes) into the same pipeline
stage.

\textbf{Common failure modes.} Local state committed to the repo; API keys in
\texttt{.tfvars}; schema change and infrastructure change validated separately and shipped
incompatibly; down-migrations attempted under pressure instead of forward-only fixes.

\section{Chapter 22: Minikube Operator Lab}
\textbf{Approach.} Start Minikube, install the MongoDB Community Operator via its manifests,
apply a \texttt{MongoDBCommunity} custom resource describing a three-member replica set with
TLS, then drill the reconciliation loop: delete a pod and watch it heal, roll an upgrade, and
confirm members rejoin with health gates between steps.

\textbf{Key commands.} \texttt{minikube start}; \texttt{kubectl apply -f
mongodbcommunity.yaml}; \texttt{kubectl get pods -w}; \texttt{kubectl delete pod <member-1>} and
observe; \texttt{kubectl describe mongodbcommunity} for status conditions.

\textbf{Common failure modes.} Expecting the operator to choose shard keys or sizing (it
reconciles; it does not design); PVCs left behind after \texttt{minikube delete} consuming disk;
TLS secrets created in the wrong namespace; judging readiness by pod Running state instead of
replica-set health.

\section{Chapter 23: The PITR Restore Drill}
\textbf{Approach.} Pick a target timestamp, restore the preceding snapshot to a \emph{new}
cluster, replay retained oplog entries to the exact point, run the validation queries (counts,
checksums, application smoke tests), and record measured RTO and RPO against their targets.
Finish with the runbook update the drill revealed.

\textbf{Common failure modes.} Restoring onto the production cluster; assuming PITR granularity
equals snapshot interval rather than oplog retention; drills executed only by the runbook author;
backup job failures unmonitored until restore day.

\section{Chapter 24: The Timed Evaluation and Gap Sprints}
\textbf{Approach.} Sit the chapter's timed self-evaluation under exam conditions, score it
against the answer rubric, and convert every missed objective into a one-day gap sprint: rebuild
the lab, re-derive the concept from the chapter, and re-test within 48 hours. Close with the
Milestone 6 capstone: the full GitOps platform deployment defended with evidence.

\textbf{Common failure modes.} Studying the evaluation instead of sitting it; gap ``review''
without a rebuilt lab (recognition is not recall); the capstone demoed without its runbook,
checklist evidence, or cost teardown.

\begin{tipbox}
For every lab, the transcript is the deliverable. Save the commands, the expected output, and the
one thing that surprised you; that habit is what Chapters 23 and 24 later grade as
production-grade judgment.
\end{tipbox}
